%% ═══════════════════════════════════════════════════════════════════════════════
\section{Two-Network Transfer Encoding with Zonotope Bounds}
\label{sec:no_n2_xp}
%% ═══════════════════════════════════════════════════════════════════════════════

The full transfer formulation (Section~\ref{sec:transfer}) encodes three forward passes---$N_1(x)$, $N_2(x)$, and $N_2(\xp)$---each introducing ReLU binary variables. Since $N_2(\xp)$ accounts for roughly one-third of all binaries, we propose a relaxed encoding that replaces the explicit $N_2(\xp)$ pass with interval-bounded output variables derived from zonotope propagation, reducing the MIP to two encoded networks while preserving soundness.

\subsection{Encoding Overview}

The MIP encodes exactly two forward passes:

\begin{center}
\begin{tikzpicture}[
  box/.style={draw,rounded corners,minimum width=2.4cm,minimum height=0.9cm,
               text centered,fill=blue!10},
  ibox/.style={draw,rounded corners,minimum width=2.4cm,minimum height=0.9cm,
               text centered,fill=orange!15,dashed},
  arr/.style={-Stealth,thick}]
  \node[box] (n1c) {$N_1(x)$};
  \node[box,right=2cm of n1c] (n2c) {$N_2(x)$};
  \node[ibox,below=1.2cm of n2c] (n2p) {$N_2(\xp)$ bounds};
  \draw[arr] (n1c) -- node[above,font=\small]{same $x$} (n2c);
  \draw[arr,dashed] (n2c) -- node[right,font=\small]{zonotope} (n2p);
\end{tikzpicture}
\end{center}

\begin{enumerate}
  \item \textbf{$N_1(x)$} is encoded with full ReLU big-$M$ constraints (binary variables for every split neuron).
  \item \textbf{$N_2(x)$} is encoded identically, sharing the same clean input variables $x$.
  \item \textbf{$N_2(\xp)$} is \emph{not} encoded as a network pass. Instead, its output logits are represented by interval-bounded variables linked to $N_2(x)$'s outputs via pre-computed perturbation-difference bounds.
\end{enumerate}

\subsection{Zonotope Perturbation-Difference Propagation}
\label{sec:no_n2_xp:zonotope}

To bound $N_2(\xp) - N_2(x)$ at the output layer without encoding $N_2(\xp)$, we propagate the perturbation difference through $N_2$ using a zonotope abstraction.

\paragraph{Initialization.}
At the input layer, the difference $\xp - x = e$ lies in $[-\varepsilon, \varepsilon]^n$. We represent this as a zonotope:
\[
  d^{[0]} \;=\; 0 + \sum_{i=1}^{n} \varepsilon \cdot g_i, \qquad g_i = e_i,
\]
where $e_i$ is the $i$-th standard basis vector and each noise symbol $\xi_i \in [-1, 1]$.

\paragraph{Linear layers.}
For a linear layer with weight matrix $W$ and bias $b$, the difference propagates exactly. Let $W_1, W_2$ be the weight matrices of $N_1$ and $N_2$ at this layer, and $\Delta W = W_2 - W_1$, $\Delta b = b_2 - b_1$. The zonotope center and generators update as:
\[
  c^{[l+1]} = W_2\, c^{[l]} + \Delta W\, \mu^{[l]}_{N_1} + \Delta b, \qquad
  G^{[l+1]} = W_2\, G^{[l]},
\]
where $\mu^{[l]}_{N_1}$ is the midpoint of $N_1$'s activation interval at layer $l$, accounting for the weight difference between the two networks.

\paragraph{Convolutional layers.}
When \texttt{zonotope\_conv} is enabled, the zonotope is propagated directly through convolutional layers using the same linear update rule applied to the 4D filter tensors, activating the zonotope at the first convolution rather than at the flatten layer.

\paragraph{ReLU layers (DeepZ).}
At a ReLU, the zonotope is updated using the DeepZ abstraction. For each neuron $j$:
\begin{itemize}
  \item If the pre-activation interval $[\ell_j, u_j]$ satisfies $\ell_j \ge 0$ (stable-active): the zonotope passes unchanged.
  \item If $u_j \le 0$ (stable-inactive): the center and generator row are zeroed.
  \item Otherwise (split): the standard DeepZ $\lambda$-relaxation introduces one fresh noise symbol per split neuron, with slope $\lambda_j = u_j / (u_j - \ell_j)$.
\end{itemize}

\paragraph{Refined ReLU case analysis.}
When \texttt{refined\_relu\_zonotope} is enabled, tighter sub-case bounds are used when one network's neuron is stable-active while the other is split. Instead of applying the generic DeepZ relaxation, the stable-active case constrains the difference to a narrower interval, reducing the generator count and bound width at no extra cost.

\paragraph{Output bounds.}
After propagating through all layers, the zonotope yields per-class bounds on the output difference:
\[
  p^{\mathrm{lo}}_k \;\le\; \hat{y}^{N_2}_k(\xp) - \hat{y}^{N_2}_k(x) \;\le\; p^{\mathrm{up}}_k, \qquad \forall\, k \in C,
\]
computed as $p^{\mathrm{up}}_k = c_k + \|G_{k,:}\|_1$ and $p^{\mathrm{lo}}_k = c_k - \|G_{k,:}\|_1$.

\subsection{Interval-Bounded Output Variables}

Using the zonotope-derived bounds, we introduce free variables $\hat{y}^{N_2}_k(\xp)$ for each output class $k$ and constrain them to lie in the interval around $N_2(x)$'s encoded outputs:
\[
  \hat{y}^{N_2}_k(x) + p^{\mathrm{lo}}_k
  \;\le\;
  \hat{y}^{N_2}_k(\xp)
  \;\le\;
  \hat{y}^{N_2}_k(x) + p^{\mathrm{up}}_k,
  \qquad \forall\, k \in C.
\]

These $|C|$ bounded variables replace the entire $N_2(\xp)$ encoding---all its ReLU binary variables, big-$M$ constraints, and intermediate layer variables are eliminated.

\subsection{Transfer MIP Formulation}

With the interval-bounded $N_2(\xp)$ outputs, the transfer MIP becomes:
\begin{equation}
  \label{eq:transfer_no_n2xp}
  \dD \;=\; \max_{x}\; \bigl(\conf{N_2}{x}{c} - \conf{N_1}{x}{c}\bigr)
  \quad\text{s.t.}\quad
  \begin{cases}
    \conf{N_1}{x}{c} \;\ge\; \dO + 10^{-3} & \text{(T1: $N_1$ certifiably robust)} \\[2pt]
    \conf{N_2}{x}{c} - \conf{N_1}{x}{c} \;\ge\; 0 & \text{(T2: $N_2$ at least as confident)} \\[2pt]
    \conf{N_2}{\xp}{c} \;\ge\; 10^{-3} & \text{(T3: $N_2$ fooled under perturbation)} \\[2pt]
    \hat{y}^{N_2}_k(\xp) \in \bigl[\hat{y}^{N_2}_k(x) + p^{\mathrm{lo}}_k,\; \hat{y}^{N_2}_k(x) + p^{\mathrm{up}}_k\bigr] & \text{(zonotope bounds)}
  \end{cases}
\end{equation}

where $\conf{N_2}{\xp}{c}$ is computed from the interval-bounded output variables, and T1 and T2 use the exact MIP-encoded confidence margins of $N_1(x)$ and $N_2(x)$.

\subsection{Soundness}

The encoding is sound because the zonotope bounds are an over-approximation: every reachable output of $N_2(\xp)$ satisfies the interval constraints. Therefore, the optimal $\dD$ from~\eqref{eq:transfer_no_n2xp} is an \emph{upper bound} on the true transfer delta-diff. Any input certified robust under this relaxed MIP is also certified under the exact three-network encoding.

\begin{remark}
  Since the zonotope over-approximates $N_2(\xp)$'s reachable outputs, the MIP may find a feasible solution where the interval-bounded variables take values that no actual $\xp$ achieves. This means $\dD$ from~\eqref{eq:transfer_no_n2xp} satisfies $\dD \ge \dD^{\mathrm{exact}}$, yielding a conservative (but sound) robustness certificate.
\end{remark}
